 %% ============================================================================
  % !Mode:: "TeX:UK:UTF-8"
  % !TEX program = XeLaTeX
  % !BIB program = biblatex
  % -----------------------------------------------------------------------------
\documentclass[notitlepage,british]{scrartcl}

\usepackage{manual} % the things usually needed in scientific texts
\usepackage{alchemist}% alchemist symbols of unicode

\hypersetup{final=true,
            pdfauthor={Engelbert Buxbaum},
            colorlinks=true,
            citecolor=Blue,%bibliography, dark blue
            linkcolor=Blue,%internal links, dark blue
            urlcolor=DodgerBlue,%internet, middle blue
            allbordercolors=White,
            pdftitle={Alchemist symbols in Unicode}}

\title{Alchemist symbols in Unicode}
\author{Engelbert Buxbaum}
\date{\today}

\addbibresource{alchemist-ref.bib}

\begin{document}
%\let\male\wasymale
%\let\female\wasyfemale

\maketitle
\begin{refsection}

The alchemist-style  makes accessible alchemistic and astrological symbols in Unicode. It uses the Gnu Unifont \href{https://unifoundry.com/unifont/}{https://unifoundry.com/unifont/}, which has all required symbols. Unfortunately, the font is split into an upper and lower part. The Symbola-font (\href{https://fontlibrary.org/de/font/symbola}{https://fontlibrary.org/de/font/symbola}) and Quivira (\href{https://www.fontspace.com/quivira-font-f13271}{https://www.fontspace.com/quivira-font-f13271}) may also be suitable, but are incomplete. If you want to use them, set AlchemistA and AlchemistB to the same font. Quivira has a more handwritten look that may fit very well to the ancient character of the symbols.

\section{Concepts}

The symbols used by alchemists served as \Foreign{aide memoir} for the researchers themself and for communication, even across language barriers. At the same time, however, their meaning was hidden except to a small group of specially trained practitioners.

Greek cosmology is described in the \foreignlanguage{greek}{Τίμαιος} (Tímaios) by \foreignlanguage{greek}{Πλάτων} \textsc{Plátōn} (ca. \num{428}--\num{348} BC) \parencite{Pla-40}, but goes back to earlier philosophers (\foreignlanguage{greek}{Ἐμπεδοκλῆς} Empedoklḗs ca. \num{490}--\num{430} BC, \foreignlanguage{greek}{Δημόκριτος} Dēmókritos ca. \num{460}--\num{370} BC). Assume a cube of some material, say, iron. This cube could be cut in halves, those halves into halves again, and so on. However, this is not possible \Foreign{ad infinitum}, but at some point one would reach a limit, where further division is no longer possible. These smallest possible particles are called atoms (from Greek \foreignlanguage{greek}{ἄτομος} átomos = undividable). There are \num{4} kinds of atoms, shaped like those \textsc{Platon}ic bodies constructed from triangles: tetrahedron, octahedron, icosahedron and cube. These different atomic shapes correspond to the 4 elements: the atoms of fire are tetrahedrons (prickly) , soil cubes (tessellates Euclidean space and causes the solidity of the soil), water icosahedras (flows out of the hands like little balls) and air octahedrons (smooth, can hardly be felt). Each of those elements embodies a particular combination of the primary qualities hot -- cold and wet -- dry: air is hot and wet, fire hot and dry, soil dry and cold and water wet and cold. Air is gaseous, water liquid and soil solid, fire is consuming. Fire represents the male, water the female principle, air and soil are transition elements.

Each of the basic elements is also associated with one of the 4 bodily fluids that must be in balance to avoid disease, they also correspond to temperaments: blood -- air -- sanguinic, black gall -- soil -- melancholic, yellow gall -- fire -- choleric and slime -- water -- phlegmatic (humoral pathology, \foreignlanguage{greek}{Ἱπποκράτης} Hippokrates \num{460}--\num{370} BC, \foreignlanguage{greek}{Γαληνός} \textsc{Galenos}, \num{129}--\num{199} AD, \foreignlanguage{arabic}{أبو علي الحسين بن عبد الله ابن سينا} Abū Alī al-Husain ibn Abd Allāh ibn Sīnā (Avicenna) \num{980}--\num{1037} AD).

The fifth platonic body -- dodecaeder -- is not composed of triangles, is considered the most noble and represents the cosmos or the spirit.

Alchemy had three goals:
\begin{itemize}
  \item{the preparation of chemicals, often in the context of practical purposes (\Foreign{e.g.}, metals and dyes). This is now handled by chemistry.}
  \item{the preparation of medicines that can cure diseases. This is now dealt with by pharmacology.}
  \item{the transmutation of base metals (like lead) into noble metals like silver and gold (Gr. \foreignlanguage{greek}{χρυσοποιία}, chrysopoeia = gold making). Today we understand that this requires the transformation of elements (nuclei with different number of protons), which is not possible by chemistry. Nuclear physics can do so in particle accelerators by bombarding bismuth with high energy atoms to kick out \num{4} protons and \num{6}--\num{9} neutrons from the \isotope[209][83]{Bi}-nuclei to produce a variety of \isotope[][79]{Au}-isotopes, of which only \isotope[197][79]{Au} is not radioactive \parencite{Ale-81}. Also, the process requires vast amounts of energy (worth about US\$ \num{5000} per \si{h}, for \SI{1}{d} to produce a few thousand gold atoms) and is thus not cost effective. For true alchemists, however, transmutation was only a symbol for obtaining spiritual enlightenment by liberating ones essence from the worldly personality.}
\end{itemize}

Later authors added the quintessence, an eternal substance that forms the basis of the other elements. \Name{Philippus Theophrastus Aureolus Bombast von Hohenheim (Paracelsus)} \num{1493}--\num{1541} \parencite{Mic-21} replaced the 4 elements from antiquity with the \Foreign{tria prima} (original trias) mercury (volatile), sulfur (burning) and salt (stabilising). He didn't mean the substances we take from the lab shelf, but their ``philosophical'' equivalent. \Name{Paracelsus} based medicine on observation rather than philosophical book wisdom and explored the pharmaceutical use of minerals, complementing the plants suggested by \Name{Galenos}. He introduced the concept that substances toxic in higher doses may be curative in low (\Foreign{dosis facit venenum}, the dose makes the poison) and is the father of modern toxicology. In particular, he experimented with mercury against the ``French disease'' (syphilis), but the results were, at best, variable. It is thus not surprising that \Name{Paracelsus} died from chronic exposure to mercury vapours as attested by the \chemical{Hg}-concentration in his bones.

\begin{landscape}
  \tablehead{
        \toprule
        Hex    & Decimal       & Char                   & Name         & Description \\
        \midrule }
  \tabletail{\midrule \multicolumn{5}{r}{continued on next page} \\}
  \tablelasttail{\bottomrule}
  \begin{supertabular}{rrrll}
    1F747 & \num{128839}  & \AlchemistSpirit       & Spirit       &  \\
    1F700 & \num{128768}  & \AlchemistQuintessence & Quintessence & aether \\
    1F701 & \num{128769}  & \AlchemistAir          & Air          &  \\
    1F702 & \num{128770}  & \AlchemistFire         & Fire         &  \\
    1F703 & \num{128771}  & \AlchemistSoil         & Soil         &  \\
    1F704 & \num{128772}  & \AlchemistWater        & Water        &  \\
    1F70E & \num{128782}  & \AlchemistPhilosophersSulphur & Philosopher's Sulphur & \\
     2721 &  \num{10017}  & \AlchemistMateriaPrima        & \Foreign{Materia prima} & \\
  \end{supertabular}
\end{landscape}

\section{Compounds}

\begin{landscape}
  \tablehead{
        \toprule
        Hex    & Decimal       & Char                   & Name         & Description \\
        \midrule }
  \tabletail{\midrule \multicolumn{5}{r}{continued on next page} \\}
  \tablelasttail{\bottomrule}
  \begin{supertabular}{rrrll}
    1F705 & \num{128773} & \AlchemistAquaFortis          & Aqua fortis            & Nitric acid \chemical{HNO_3} \\
    1F706 & \num{128774} & \AlchemistAquaRegiaA          & Aqua regia 1           & 1 part nitric, 3 parts hydrochloric acid \chemical{NOCl} \\
    1F707 & \num{128775} & \AlchemistAquaRegiaB          & Aqua regia 2           &                              \\
    1F708 & \num{128776} & \AlchemistAquaVitaeA          & Aqua vitae 1           & Ethanol, distilled spirit \chemical{CH_3\single CH_2\single OH}   \\
    1F709 & \num{128777} & \AlchemistAquaVitaeB          & Aqua vitae 2           &  \\
    1F70A & \num{128778} & \AlchemistVinegar             & Vinegar                & acetic acid \chemical{CH_3\single COOH} \\
    1F70B & \num{128779} & \AlchemistDistilledVinegarA   & distilled vinegar 1    & \\
    1F70C & \num{128780} & \AlchemistDistilledVinegarB   & distilled vinegar 2    &  \\
    1F710 & \num{128784} & \AlchemistSublimateOfMercuryA & Sublimate of mercury 1 & mercury(II)chloride \chemical{HgCl_2} \\
    1F711 & \num{128785} & \AlchemistSublimateOfMercuryB & Sublimate of mercury 2 &  \\
    1F712 & \num{128786} & \AlchemistSublimateOfMercuryC & Sublimate of mercury 3 &  \\
    1F713 & \num{128787} & \AlchemistCinnabar            & Cinnabar               & mercury(II)sulphide \chemical{HgS} \\
    1F714 & \num{128788} & \AlchemistSalt                & Salt                   & sodium chloride \chemical{NaCl} \\
    1F715 & \num{128789} & \AlchemistNitre               & Nitre                  & salpeter, \chemical{KNO_3} \\
    1F716 & \num{128790} & \AlchemistVitriolA            & Vitriol 1              & sulphuric acid \chemical{H_2SO_4} \\
    1F717 & \num{128791} & \AlchemistVitriolB            & Vitriol 2              &  \\
    1F718 & \num{128792} & \AlchemistRockSaltA           & Rock salt 1            & \textit{sal gemmae}, impure \chemical{NaCl} \\
    1F713 & \num{128793} & \AlchemistRockSaltB           & Rock salt 2            &  \\
    1F71C & \num{128796} & \AlchemistIronOreA            & Iron ore 1             &  \\
    1F71D & \num{128797} & \AlchemistIronOreB            & Iron ore 2             &  \\
    1F71E & \num{128798} & \AlchemistCrocusOfIron        & Crocus of iron         & \textit{crocus martis}, rust \\
    1F720 & \num{128800} & \AlchemistCopperOre           & Copper ore             &  \\
    1F723 & \num{128803} & \AlchemistCrocusOfCopperA     & Crocus of copper 1     & copper(II)oxyde \chemical{CuO} \\
    1F724 & \num{128804} & \AlchemistCrocusOfCopperB     & Crocus of copper 2     &  \\
    1F721 & \num{128801} & \AlchemistIronCopperOre       & Iron copper ore        & chalcopyrite \chemical{CuFeS_2} \\
    1F71B & \num{128795} & \AlchemistSublimateOfCopperA  & Sublimate of copper 1  & copper(I)chloride \chemical{CuCl} \\
    1F71C & \num{128796} & \AlchemistSublimateOfCopperB  & Sublimate of copper 2  &  \\
    1F725 & \num{128805} & \AlchemistCopperAntimonateA   & Copper antimonate 1    & \chemical{CuSb_2O_6} \\
    1F726 & \num{128806} & \AlchemistCopperAntimonateB   & Copper antimonate 2    & \\
    1F728 & \num{128808} & \AlchemistVerdigris           & Verdigris              & patina on copper (acetate, hydroxyde, chloride) \\
    1F729 & \num{128809} & \AlchemistTinOre              & Tin ore                & Cassiterite \chemical{SnO_2} \\
    1F72A & \num{128810} & \AlchemistLeadOre             & Lead ore               & Galena \chemical{PbS} \\
    1F72B & \num{128811} & \AlchemistAntimonyOre         & Antimony ore           & stibnite \chemical{Sb_2S_3} \\
    1F72C & \num{128812} & \AlchemistSublimateOfAntimony & Sublimate of antimony  & senarmontite \chemical{Sb_4O_6}  \\
    1F72D & \num{128813} & \AlchemistSaltOfAntimony      & Salt of antimony       & Antimony sulphate \chemical{Sb_2(SO_4)_3} \\
    1F72E & \num{128814} & \AlchemistSublimateOfSaltOfAntimony & Sublimate of salt of antimony & an antimony oxyde from decomposition of the sulphate? \\
    1F72F & \num{128815} & \AlchemistVinegarOfAntimony   & Vinegar of antimony    & may be \chemical{Sb(CH_3COO)_3}, but this is unstable \\
    1F736 & \num{128822} & \AlchemistAlkaliA             & Alkali 1               & \Foreign{sal alkali} \chemical{Na_2CO_3} or \chemical{K_2CO_3} \\
    1F737 & \num{128823} & \AlchemistAlkaliB             & Alkali 2               &  \\
    1F738 & \num{128824} & \AlchemistMarcasite           & Marcasite              & \chemical{FeS_2} (originally included iron pyrite or fools gold) \\
    1F739 & \num{128825} & \AlchemistSalAmmoniak         & Sal ammoniak           & ammonium chloride \chemical{NH_4Cl} \\
    1F73B & \num{128827} & \AlchemistRealgarA            & Realgar 1              & \chemical{As_4S_4} \\
    1F73C & \num{128828} & \AlchemistRealgarB            & Realgar 2              &  \\
    1F73D & \num{128829} & \AlchemistAuripigment         & Auripigment            & \chemical{As2S3} \\
    1F73E & \num{128830} & \AlchemistBismuthOre          & Bismuth ore            & bismuthinite \chemical{Bi_2S_3} or bismite \chemical{Bi_2O_3}?\\
    1F73F & \num{128831} & \AlchemistTartarA             & Tartar1                & wine stone,  potassium hydrogen tartrate \\
    1F740 & \num{128832} & \AlchemistTartarB             & Tartar2                &  \\
    1F741 & \num{128833} & \AlchemistQuicklime           & Quicklime              & calcium oxyde \chemical{CaO}\\
    1F742 & \num{128834} & \AlchemistBoraxA              & Borax 1                & \chemical{Na_2B_4O_7 \times 10 H_2O} \\
    1F743 & \num{128835} & \AlchemistBoraxB              & Borax 2                &  \\
    1F744 & \num{128836} & \AlchemistBoraxC              & Borax 3                &  \\
    1F745 & \num{128837} & \AlchemistAlum                & Alum                   & double salt like \chemical{KAl(SO_4)_2 \times 12  H_2O} \\
    1F753 & \num{128851} & \AlchemistLodestone           & Lodestone              & magnetite, Fe(II, III)oxyde \\
    1F758 & \num{128856} & \AlchemistPotashes            & Potashes               & \chemical{K_2CO_3} \\
    1F75C & \num{128860} & \AlchemistStratumSuperStratumA & \textit{Stratum super stratum} 1 & layer over layer \\
    1F75D & \num{128861} & \AlchemistStratumSuperStratumB & \textit{Stratum super stratum} 2 &  \\
  \end{supertabular}
\end{landscape}

\section{Mixtures}

\begin{landscape}
  \tablehead{
        \toprule
        Hex    & Decimal       & Char                   & Name         & Description \\
        \midrule }
  \tabletail{\midrule \multicolumn{5}{r}{continued on next page} \\}
  \tablelasttail{\bottomrule}
  \begin{supertabular}{rrrll}
    1F70F & \num{128863} & \AlchemistBlackSulphur        & Black sulphur          & residue left after sublimation of sulphur \\
    1F754 & \num{128852} & \AlchemistSoap                & Soap                   & sodium salt of mixed fatty acids \\
    1F755 & \num{128853} & \AlchemistUrine               & Urine                  &  \\
    1F756 & \num{128854} & \AlchemistHorseDung           & Horse dung             &  \\
    1F757 & \num{128855} & \AlchemistAshes               & Ashes                  &  \\
    1F759 & \num{128857} & \AlchemistBrick               & Brick                  &  \\
    1F75A & \num{128858} & \AlchemistPowderedBrick       & powdered brick         &  \\
    1F75B & \num{128859} & \AlchemistAmalgam             & Amalgam                & alloy of mercury \\
    1F74E & \num{128862} & \AlchemistCaputMortuum        & \textit{Caput mortuum} & slag left after a purification \\
    1F746 & \num{128838} & \AlchemistOil                 & Oil                    &  \\
    1F749 & \num{128841} & \AlchemistGum                 & Gum                    &  \\
    1F74A & \num{128842} & \AlchemistWax                 & Wax                    &  \\
    1F74C & \num{128844} & \AlchemistCalx                & Calx                   & oxydised residue, also \chemical{CaO} \\
    1F74D & \num{128845} & \AlchemistTutty               & Tutty                  & \chemical{ZnO} or \chemical{ZnCO_3} \\
  \end{supertabular}
\end{landscape}

\section{Elements}

The alchemist's idea about elements (atoms shaped like different \Name{Platon}ic bodies) was very different from ours (different number of protons in the nucleus). In the following, we list elements in our modern understanding. Even then, however, there can be misunderstandings. ``Regulus of antimony'' is quite clear, a drop of antimony left after purification. However, ``regulus of iron'' and ``regulus of copper'' also mean antimony, but made by reducing antimony oxyde with iron or copper, respectively. It was assumed that they were different because of the different synthetic route.

The symbols for iron/Mars and copper/Venus are also used for male and female.

\begin{landscape}
  \tablehead{
        \toprule
        Hex    & Decimal       & Char                   & Name         & Description \\
        \midrule }
  \tabletail{\midrule \multicolumn{5}{r}{continued on next page} \\}
  \tablelasttail{\bottomrule}
  \begin{supertabular}{rrrll}
    2646  &   \num{9798} & \AlchemistBismuth             & Bismuth                & also planet Neptune \\
    2640  &   \num{9792} & \AlchemistCopper              & Copper                 & also planet Venus \\
    1F71A & \num{128794} & \AlchemistGoldA               & Gold                   &  \\
    2609  &   \num{9737} & \AlchemistGoldB               & Gold                   & also sun   \\
    2642  &   \num{9794} & \AlchemistIron                & Iron                   & also planet Mars \\
    2644  &   \num{9796} & \AlchemistLead                & Lead                   & also planet Saturn \\
    1F730 & \num{128816} & \AlchemistRegulusOfAntimonyA  & Regulus of Antimony 1  & impure antimony \\
    1F731 & \num{128817} & \AlchemistRegulusOfAntimonyB  & Regulus of Antimony 2  &  \\
    1F71F & \num{128799} & \AlchemistRegulusOfIron       & Regulus of iron        & antimony prepared from stibnite by reduction with iron \\
    263F  &   \num{9791} & \AlchemistMercury             & Mercury                & also the planet mercury \\
    1F71B & \num{128795} & \AlchemistSilverA             & Silver 1               &  \\
    263D  &   \num{9789} & \AlchemistSilverB             & Silver 2               & also first quarter moon  \\
    263E  &   \num{9790} & \AlchemistSilverC             & Silver 3               & also last quarter moon \\
    1F70D & \num{128781} & \AlchemistSulphur             & Sulphur                &  \\
    2643  &   \num{9795} & \AlchemistTin                 & Tin                    & also planet Jupiter \\
  \end{supertabular}
\end{landscape}

\section{Processes}

As mentioned, the \Foreign{opus magnum} of gold making is really only an allegory of the journey to self-knowledge and, yes, eventually to god \parencite{Rip-71, Jun-44}:

\begin{description}
  \item[Calcination]{exposure of a sample to high, dry heat (roasting, even glowing). Many substances under these condition loose weight, turn white and become easy to grind. First (black) stage of the \Foreign{opus magnum}. The colour black represents chaos, and that what is hidden or buried, the \Foreign{materia prima} from which all other things may be obtained. Spiritually, calcination means burning off our attachments to the world: the desire for status, fame, wealth and identity. Our preconceived notions about ourself are put to the test by fire, in an existential crisis, in the \emph{dark night of the soul}.}
  \item[Dissolution]{forming a \(\rightarrow\) solution. In the \Foreign{opus magnum}, the ashes obtained during calcination are dissolved in water, the symbol of the unconscious or hidden. Spiritually, we free ourselves from our inauthentic and acquired traits. The dissolution stage involves freeing repressed emotions from traumatic events that we have pushed down into our subconsciousness (dissolution of the ego \parencite{Jun-44}), and can be a cathartic step.}
  \item[Separation]{is the 3rd stage of the \Foreign{opus magnum}, where the products of dissolution are filtered and separated. Whilst the first stages were associated with the elements fire and water, this one is associated with air. The pure essence is extracted from the mixture. Spiritually, we leave our acquired personality in form of engrained thought processes and emotional triggers to become our true, essential self. Collect all the things in you that are helpful for yourself and for others, and let go of everything else. After the turmoil of the first steps, we can now enjoy the stillness of being removed from all that is inauthentic and artificial. }
  \item[Conjunction]{means bringing together the elements purified in the first three stages by fire, water and air in the sign of the soil. Spiritually, we bring what is left of us together to form a new, authentic personality. The conflict of dualities like body and soul, spirit and matter, conscious and unconscious are resolved. Feminine qualities like emotion and intuition join male like intellect and logic.  }
  \item[Fermentation]{exposure to microorganisms for longer time to achieve a chemical reaction and remove what is no longer needed. }
  \item[Putrefaction]{leaving a sample undisturbed, often after adding the desired material (seed). Both fermentation and putrefaction form the 5th stage of the \Foreign{opus magnum}: the removal of the old, unauthentic self. This process can be painful, but we cannot see dawn before we have gone through the darkest of nights. }
  \item[Distillation]{process of vaporising and then condensing a substance. The different boiling points of the components of a mixture are used for purification of the essence. In spiritual alchemy, our core identity is freed from any inferior elements. The ego is no longer dominating, so that the soul can be heard. }
  \item[Coagulation]{old term for crystallisation. The result is the philosophers stone. Spiritually, the self comes together in a healing process. This is called \emph{rubedo}. This word also refers to the state achieved by the mage where he is fully attuned to the \Foreign{anima mundi}, the fountainhead of creation and source of the \Foreign{prima materia}. }
\end{description}

Apart from the processes involved in the \Foreign{opus magnum}, the following methods were used by alchemists:
\begin{description}
  \item[Ceration]{allowing a hard, dry substance to absorb water while it is heated (imbibition) for the purpose of softening it (from Latin \Foreign{cera} = wax). }
  \item[Digestion]{heating a sample in an open vessel for long times, but without boiling it.}
  \item[Fixation]{turning a volatile compound into a solid that is not affected by fire}
  \item[Multiplication]{increasing the potency of an elixir or philosophers stone, often by repeating the process by which it was originally obtained}
  \item[Precipitation]{letting a solid form from a solution and collecting it. }
  \item[Projection]{process of transmuting base metals to gold using the philosopher's stone}
  \item[Purify]{increasing the concentration of a particular substance, if possible close to \SI{100}{\%}. }
  \item[Regulus formation]{getting the pure form of a metal (especially antimony) out of an ore. }
  \item[Solution]{one component (solute) is homogenously dispersed in a second component (solvent, usually, but not necessarily, a liquid). }
  \item[Sublimation]{heating a substance to turn it from the solid directly into the gaseous state and back, without forming a liquid first. The sublimate is called ``flower''.}
  \item[Tincture]{extract in alcohol.}
\end{description}

\begin{landscape}
  \tablehead{
        \toprule
        Hex    & Decimal       & Char                   & Name         & Description \\
        \midrule }
  \tabletail{\midrule \multicolumn{5}{r}{continued on next page} \\}
  \tablelasttail{\bottomrule}
  \begin{supertabular}{rrrll}
    2648  & \num{9800}   & \AlchemistCalcination         & Calcination            & also aries         \\
    2650  & \num{9802}   & \AlchemistCeration            & Ceration               & also Sagittarius   \\
    2649  & \num{9801}   & \AlchemistCongelation         & Congelation            & also Taurus        \\
    264C  & \num{9804}   & \AlchemistDigestion           & Digestion              & also Leo           \\
    1F761 & \num{128865} & \AlchemistDissolveA           & Dissolve 1             &                    \\
    1F762 & \num{128866} & \AlchemistDissolveB           & Dissolve 2             &                    \\
    1F760 & \num{128864} & \AlchemistDistillationA       & Distillation 1         &                    \\
    264D  & \num{9805}   & \AlchemistDistillationB       & Distillation 2         & also Virgo         \\
    2651  & \num{9803}   & \AlchemistFermentation        & Fermentation           & also Capricorn     \\
    264A  & \num{9802}   & \AlchemistFixation            & Fixation               & also Gemini        \\
    2652  & \num{9804}   & \AlchemistMultiplication      & Multiplication         & also Aquarius      \\
    1F75F & \num{128863} & \AlchemistPrecipitation       & Precipitation          &                    \\
    2653  & \num{9805}   & \AlchemistProjection          & Projection             & also Pisces        \\
    1F763 & \num{128867} & \AlchemistPurify              & Purify                 &                    \\
    1F764 & \num{128868} & \AlchemistPutrefaction        & Putrefaction           &                    \\
    1F732 & \num{128818} & \AlchemistRegulusA            & Regulus 1              & smelting of metallic ore \\
    1F733 & \num{128819} & \AlchemistRegulusB            & Regulus 2              &  \\
    1F734 & \num{128820} & \AlchemistRegulusC            & Regulus 3              &  \\
    1F735 & \num{128821} & \AlchemistRegulusD            & Regulus 4              &  \\
    264F  & \num{9807}   & \AlchemistSeparation          & Separation             & also Scorpio       \\
    264B  & \num{9803}   & \AlchemistSolution            & Solution               & also Cancer        \\
    1F75E & \num{128862} & \AlchemistSublimationA        & Sublimation 1          &                    \\
    264E  & \num{9806}   & \AlchemistSublimationB        & Sublimation 2          & also Libra         \\
    1F748 & \num{128840} & \AlchemistTincture            & Tincture               & alcoholic extract  \\
  \end{supertabular}
\end{landscape}

\section{Tools}

\begin{description}
  \item[Alembic]{from Greek \foreignlanguage{greek}{ἄμβιξ} ambix = beaker over Arabic \foreignlanguage{arabic}{الإنبيق} al-inbīq is a distillation head connecting two vessels. In the fist (cucurbit from \foreignlanguage{greek}{βῖκος} bîkos) a sample was heated, the second (receiver, from \foreignlanguage{greek}{φιάλη} phialē over Arabic \foreignlanguage{arabic}{قَابِلَة} qābila) received the vapours condensed in a connecting tube (Greek \foreignlanguage{greek}{σωλήν} sōlēn). Its invention is attributed to \Name{Maria Prophetissa} of Alexandria in the 1st century AD (variously also called \Name{Maria the Jewess} or \Name{Maria the Copt}). The bridge of modern chemical glassware is the closest equivalent.}
  \item[Athanor]{furnance}
  \item[Caduceus]{is the staff of Hermes (Gr.) = Mercury (Lat., god of traders and thieves) with two snakes and wings. It should not be mixed up with the staff of Asclepios (Son of Apollo) with one snake and no wings, which is a symbol of medicine. The astrological/alchemists sign of Mercury is an older version of the caduceus, with the snakes protruding from the staff.}
  \item[Crucible]{a heat resistant ceramic vessel used to heat substances in fire. }
  \item[Balneum Mariae ]{hot water bath invented by \Name{Maria Prophetissa} and still used to keep food warm in refectories and the like (\Foreign{bain de Marie}). It consists of an outer vessel half-filled with a hot fluid (usually water) and an inner vessel that is filled with the material to be kept warm and immersed in the fluid of the outer vessel. The material in the inner vessel is kept at a constant temperature, without hot or cold spots.}
  \item[Retort]{(from Lat. \Foreign{retortus} = curved backward) glass vessel used by alchemists. Heating it was an art in itself, as the soda-lime glasses available at the time were liable to crack when heated unevenly. In the modern chemical industry, vessels of any shape and material used for pyrolysis are still called retort (production of shale oil, charcoal, recovery of \chemical{Hg} in gold mining). In the laboratory, the \Name{Liebig} condenser has largely replaced the retort. }
  \item[Scepter of Jove]{or bident, a pitchfork-like instrument with two prongs, was used in Roman time for the consecration of a place struck by lightning. It is probably a symbol for Jupiter's lightning bolt. The name may be derived from the young sheep sacrificed in the ritual (so young as to have only two teeth). The bident is also the symbol for Hades (Pluto), the god of the underworld. }
  \item[Staff of Asclepios]{Staff with a snake (allegedly a Aesculapian adder \Foreign{Zamenis longissimus} (\Name{Laurenti}, 1768), Colubridae)  wound around it. Originally, the ``snake'' is probably a Guinea worm (\Foreign{Dracunculus medinensis} L., Dracunculidae), which already in ancient Egypt was removed from affected limbs by slowly pulling it out and winding it around a piece of wood. }
  \item[Trident ]{is similar to the bident, but has three prongs. It is the symbol of Poseidon (Neptune), god of the seas. }
  \item[Starred Trident ]{}
  \item[Vapour bath ]{}
\end{description}

It is noteworthy that there were female alchemists at a time when females were generally limited to ``kitchen, children and church'', and that the work of these female alchemists was highly respected by their male colleagues. \Name{Maria Prophetissa} and \Name{Cleopatra the Alchemist} (Alexandria, ca. 3rd century AD, not related to the Ptolemaic queens) may be the most well-known of them.

\begin{landscape}
  \tablehead{
        \toprule
        Hex    & Decimal       & Char                   & Name         & Description \\
        \midrule }
  \tabletail{\midrule \multicolumn{5}{r}{continued on next page} \\}
  \tablelasttail{\bottomrule}
  \begin{supertabular}{rrrll}
    1F76A & \num{128874} & \AlchemistAlembicA          &  Alembic 1        & distillation head  \\
     2697 &   \num{9879} & \AlchemistAlembicB          &  Alembic 2        &   \\
     2695 &   \num{9877} & \AlchemistAsclepius         &  Staff of Asclepius & symbol of medicine \\
    1F750 & \num{128848} & \AlchemistCaduceusA         &  Caduceus 1       & staff of Hermes = Mercury (with two snakes)  \\
     2624 &   \num{9764} & \AlchemistCaduceusB         &  Caduceus 2       &   \\
    1F765 & \num{128869} & \AlchemistCrucibleA         &  Crucible 1       & ceramic, heat-resistant vessel  \\
    1F766 & \num{128870} & \AlchemistCrucibleB         &  Crucible 2       &   \\
    1F767 & \num{128871} & \AlchemistCrucibleC         &  Crucible 3       &   \\
    1F768 & \num{128872} & \AlchemistCrucibleD         &  Crucible 4       &   \\
    1F769 & \num{128873} & \AlchemistCrucibleE         &  Crucible 5       &   \\
    1F76B & \num{128875} & \AlchemistBalneumMariae     &  Balneum Mariae   & water bath  \\
    1F76D & \num{128877} & \AlchemistRetort            &  Retort           & glass vessel  \\
    1F74F & \num{128879} & \AlchemistScepterOfJove     &  Scepter of Jove  &   \\
    1F751 & \num{128849} & \AlchemistTrident           &  Trident          &   \\
    1F752 & \num{128850} & \AlchemistStarredTrident    &  Starred Trident  &   \\
    1F76C & \num{128876} & \AlchemistVapourBath        &  Vapour bath      &   \\
  \end{supertabular}
\end{landscape}


\section{Astronomy and astrology}

\subsection{Astrology}

\begin{description}
  \item[Regulus]{brightest star system (apparent magnitude of +1.35) in the constellation Leo (\textalpha\ Leonis). Actually consists of 4 stars, Regulus A is a blue-white main-sequence star, its companion has not been observed directly but is probably a white dwarf. Regulus B and C are main sequence stars.}
  \item[Ascending node]{if the plane defined by the orbit of a celestial body has an angle (inclination) \(\neq 0\) to a reference plane, then the orbit intersects the reference plane in two points. The one crossed by the north-moving body is called the ascending node (Latin \Foreign{caput draconis} = dragon's head or Greek \foreignlanguage{greek}{αναβιβάζων} anabibazon). In astrology, these terms refer to the crossings of the orbit of the moon with the apparent orbit of the sun across the sky.}
  \item[Descending note]{The node of the south-moving body (\Foreign{cauda draconis} = dragon's tail or \foreignlanguage{greek}{καταβιβάζων} catabibazon). }
  \item[Conjunction]{is the apparent meeting of two celestial bodies as seen from earth (angle \(\pm 10°\)). For example, a conjunction of sun and moon is visible as solar eclipse. ``Great conjugation'' we call the meeting of Jupiter and Saturn, its occurance \num{7} BC may have been the biblical star of Bethlehem (Mt 2\textsubscript{2}).}
  \item[Opposition]{is the situation where two celestial bodies appear (from the earth) to have an angle of \ang{180}. Usually, one of the bodies is the sun, the other then is visible during the entire night. A near opposition of moon and sun is called full moon, an exact opposition lunar eclipse. Planets have their smallest distance from the earth during their opposition.  }
  \item[Sextile]{two celestial bodies appear to have an angle of \ang{60}.}
  \item[Semisextile]{two celestial bodies appear to have an angle of \ang{30}.}
  \item[Quincunx]{two celestial bodies appear to have an angle of \ang{150}.}
  \item[Sesquiquadrate]{two celestial bodies appear to have an angle of \ang{135}.}
  \item[Lot of fortune]{ or lucky point (Lat. \Foreign{Pars Fortunae}), its exact meaning varies between ancient, arabic and western astrologers. }
  \item[Occultation]{occurs when one object is hidden from the observer by a second moving between them. If the second object is too small to hide the first completely, we call this \emph{transit}. }
  \item[Lunar eclipse]{moon becomes invisible because of an exact opposition with the sun (moon enters the shadow cast by the earth). }
\end{description}

\begin{landscape}
  \tablehead{
        \toprule
        Hex    & Decimal       & Char                   & Name         & Description \\
        \midrule }
  \tabletail{\midrule \multicolumn{5}{r}{continued on next page} \\}
  \tablelasttail{\bottomrule}
  \begin{supertabular}{rrrll}
    1F732 & \num{128818}  & \AlchemistRegulusA           & Regulus 1              & one of the brightest stars, also the process of smelting \\
    1F733 & \num{128819}  & \AlchemistRegulusB           & Regulus 2              &   \\
    1F734 & \num{128820}  & \AlchemistRegulusC           & Regulus 3              &   \\
    1F735 & \num{128821}  & \AlchemistRegulusD           & Regulus 4              &   \\
    260A  & \num{9738}    & \AlchemistAscendingNode      & AscendingNode          &   \\
    260B  & \num{9739}    & \AlchemistDescendingNode     & DescendingNode         &   \\
    260C  & \num{9740}    & \AlchemistConjunction        & Conjunction            &   \\
    260D  & \num{9741}    & \AlchemistOpposition         & Opposition             &   \\
    26B9  & \num{9913}    & \AlchemistSextile            & Sextile                &   \\
    26BA  & \num{9914}    & \AlchemistSemisextile        & Semisextile            &   \\
    26BB  & \num{9915}    & \AlchemistQuincunx           & Quincunx               &   \\
    26BC  & \num{9916}    & \AlchemistSesquiquadrate     & Sesquiquadrate         &   \\
    1F774 & \num{128884}  & \AlchemistLotOfFortune       & LotOfFortune           &   \\
    1F775 & \num{128885}  & \AlchemistOccultation        & Occultation            & hiding one celestial body by another  \\
    1F776 & \num{128886}  & \AlchemistLunarEclipse       & LunarEclipse           &   \\
  \end{supertabular}
\end{landscape}

\subsubsection{Star signs}

\begin{landscape}
  \tablehead{
        \toprule
        Hex    & Decimal       & Char                   & Name         & Description \\
        \midrule }
  \tabletail{\midrule \multicolumn{5}{r}{continued on next page} \\}
  \tablelasttail{\bottomrule}
  \begin{supertabular}{rrrll}
    2648  & \num{9800}    & \AlchemistAries              & Aries                  & also calcination    \\
    2649  & \num{9801}    & \AlchemistTaurus             & Taurus                 & also congelation    \\
    264A  & \num{9802}    & \AlchemistGemini             & Gemini                 & also fixation       \\
    264B  & \num{9803}    & \AlchemistCancer             & Cancer                 & also solution       \\
    264C  & \num{9804}    & \AlchemistLeo                & Leo                    & also digestion      \\
    264D  & \num{9805}    & \AlchemistVirgo              & Virgo                  & also distillation   \\
    264E  & \num{9806}    & \AlchemistLibra              & Libra                  & also sublimation    \\
    264F  & \num{9807}    & \AlchemistScorpio            & Scorpio                & also separation     \\
    2650  & \num{9808}    & \AlchemistSagittarius        & Sagittarius            & also ceration       \\
    2651  & \num{9809}    & \AlchemistCapricorn          & Capricorn              & also fermentation   \\
    2652  & \num{9810}    & \AlchemistAquarius           & Aquarius               & also multiplication \\
    2653  & \num{9811}    & \AlchemistPisces             & Pisces                 & also projection     \\
  \end{supertabular}
\end{landscape}


\subsection{Solar system}

\begin{landscape}
  \tablehead{
        \toprule
        Hex    & Decimal       & Char                   & Name         & Description \\
        \midrule }
  \tabletail{\midrule \multicolumn{5}{r}{continued on next page} \\}
  \tablelasttail{\bottomrule}
  \begin{supertabular}{rrrll}
    2609  &  \num{9737} & \AlchemistSun                & Sun                    & also gold                    \\
    263D  &  \num{9789} & \AlchemistFirstQuarterMoon   & First quarter moon     & also silver                  \\
    263E  &  \num{9790} & \AlchemistLastQuarterMoon    & Last quarter moon      &                              \\
    26B8  &  \num{9912} & \AlchemistBlackMoonLilith    & Black moon Lilith      & 2nd focal point of lunar orbit \\
    263F  &  \num{9791} & \AlchemistMercury            & Mercury                & also the element mercury     \\
    2640  &  \num{9792} & \AlchemistVenus              & Venus                  & also copper                  \\
    2641  &  \num{9793} & \AlchemistEarth              & Earth                  &                              \\
    2642  &  \num{9794} & \AlchemistMars               & Mars                   & also iron                    \\
    2643  &  \num{9795} & \AlchemistJupiter            & Jupiter                & also tin                     \\
    2644  &  \num{9796} & \AlchemistSaturn             & Saturn                 & also lead                    \\
    2645  &  \num{9797} & \AlchemistUranus             & Uranus                 &                              \\
    2646  &  \num{9798} & \AlchemistNeptune            & Neptune                & also bismuth                 \\
    2647  &  \num{9799} & \AlchemistPlutoA             & Pluto 1                &                              \\
    2BD3  & \num{11219} & \AlchemistPlutoB             & Pluto 2                &                              \\ 
  \end{supertabular}
\end{landscape}

\subsubsection{Dwarf planets}

\begin{landscape}
  \tablehead{
        \toprule
        Hex    & Decimal       & Char                   & Name         & Description \\
        \midrule }
  \tabletail{\midrule \multicolumn{5}{r}{continued on next page} \\}
  \tablelasttail{\bottomrule}
  \begin{supertabular}{rrrll}
    1F77B & \num{128891}  & \AlchemistHaumea             & Haumea                 &             \\
    1F77C & \num{128892}  & \AlchemistMakemake           & Makemake               &             \\
    1F77D & \num{128893}  & \AlchemistGonggong           & Gonggong               &             \\
    1F77E & \num{128894}  & \AlchemistQuaoar             & Quaoar                 &             \\
    1F77F & \num{128895}  & \AlchemistOrcus              & Orcus                  &             \\
     2BF0 &  \num{11248}  & \AlchemistErisA              & Eris 1                 &             \\
     2BF1 &  \num{11249}  & \AlchemistErisB              & Eris 2                 &             \\
     2BF2 &  \num{11250}  & \AlchemistSedna              & Sedna                  &             \\
  \end{supertabular}
\end{landscape}

\subsubsection{Asteroids}

\begin{landscape}
  \tablehead{
        \toprule
        Hex    & Decimal       & Char                   & Name         & Description \\
        \midrule }
  \tabletail{\midrule \multicolumn{5}{r}{continued on next page} \\}
  \tablelasttail{\bottomrule}
  \begin{supertabular}{rrrll}
  26B3  & \num{9907}    & \AlchemistCeres              &  Ceres                 &              \\
  26B4  & \num{9908}    & \AlchemistPallas             &  Pallas                &              \\
  26B5  & \num{9909}    & \AlchemistJuno               &  Juno                  &              \\
  26B6  & \num{9910}    & \AlchemistVesta              &  Vesta                 &              \\
  26B7  & \num{9911}    & \AlchemistChiron             &  Chiron                &              \\
  \end{supertabular}
\end{landscape}

\section{Measures}

\begin{landscape}
  \tablehead{
        \toprule
        Hex    & Decimal       & Char                   & Name         & Description \\
        \midrule }
  \tabletail{\midrule \multicolumn{5}{r}{continued on next page} \\}
  \tablelasttail{\bottomrule}
  \begin{supertabular}{rrrll}
  1F76E & \num{128878}  & \AlchemistHourA              & Hour                   &             \\
   29D6 &  \num{10710}  & \AlchemistHourB              & Hour                   &             \\
   29D7 &  \num{10711}  & \AlchemistHourC              & Hour                   &             \\
  1F76F & \num{128879}  & \AlchemistNight              & Night                  &             \\
  1F770 & \num{128880}  & \AlchemistDayNight           & DayNight               &             \\
  1F771 & \num{128881}  & \AlchemistMonth              & Month                  &             \\
  1F772 & \num{128882}  & \AlchemistHalfDram           & HalfDram               & 1 dram \(\approx\) \SI{3.6}{mL} \\
  1F773 & \num{128883}  & \AlchemistHalfOunce          & HalfOunce              & 1 ounce \(\approx\) \num{28}--\SI{31}{g} \\
  \end{supertabular}
\end{landscape}

\printbibliography[heading=subbibliography]

\end{refsection}
\end{document}
